\documentclass[11pt]{article}
\usepackage{amsmath, amssymb, graphicx, hyperref}
\usepackage{enumitem}
\setlist{nosep}
\usepackage[margin=1in]{geometry}

\title{In-Class Problem Set: Mapping Social Phenomena with US Election Data (R + GitHub \emph{or} Canvas)}
\author{}
\date{}

\begin{document}
\maketitle

\noindent \textbf{Goal.} Practice building and interpreting choropleth maps using real political data. You will join election data to state boundaries, choose appropriate color scales, compare counts vs rates, and decide when a map is the right (or wrong) chart. You may submit via \textbf{GitHub or Canvas}.

\medskip
\noindent \textbf{Dataset.} This problem set uses two sources:
\begin{itemize}
  \item \textbf{US state boundaries} from the \texttt{maps} package (\texttt{map\_data("state")}).
  \item \textbf{Built-in state data} from R's \texttt{state.x77} matrix and \texttt{state.name} vector, plus the \texttt{USArrests} dataset.
\end{itemize}

\medskip
\noindent \textbf{Key variables (state.x77).}
\begin{itemize}
  \item \texttt{Population}: estimated population (thousands, 1975 est.)
  \item \texttt{Income}: per capita income (1974)
  \item \texttt{Illiteracy}: illiteracy rate (\%, 1970)
  \item \texttt{Life.Exp}: life expectancy (years, 1969--71)
  \item \texttt{Murder}: murder rate (per 100,000, 1976)
  \item \texttt{HS.Grad}: high-school graduation rate (\%, 1970)
\end{itemize}

\noindent \textbf{Key variables (USArrests).}
\begin{itemize}
  \item \texttt{Murder}: murder arrests per 100,000
  \item \texttt{Assault}: assault arrests per 100,000
  \item \texttt{UrbanPop}: percent urban population
  \item \texttt{Rape}: rape arrests per 100,000
\end{itemize}

\medskip
\noindent \textbf{What to submit (GitHub or Canvas).}
\begin{itemize}
  \item \texttt{scripts/lab.R}
  \item \texttt{outputs/writeup.md}
  \item Required figures in \texttt{figures/}
\end{itemize}

\medskip
\noindent \textbf{Rules.}
\begin{itemize}
  \item Work inside an \textbf{R Project}.
  \item Use a \textbf{sequential, hard-coded workflow} (no user-defined functions).
  \item Save all plots using \texttt{ggsave()}.
  \item If submitting via GitHub, Git commands must be run in the \textbf{Terminal tab}.
\end{itemize}

\section*{Questions}

\begin{enumerate}

\item \textbf{Load data and prepare for joining (proof required).}

\begin{verbatim}
# load required packages
library(__________)
library(__________)
library(__________)

# get map boundary data
us <- map_data("state")

# build state attribute data
state_data <- data.frame(
  region = tolower(state.name),
  as.data.frame(state.x77)
)

# quick inspection
dim(us)
dim(state_data)
head(state_data)
\end{verbatim}

Include in your write-up:
\begin{itemize}
  \item number of rows in boundary data vs attribute data,
  \item confirmation that the \texttt{region} column matches between the two,
  \item a quick summary of at least two numeric columns.
\end{itemize}

\item \textbf{Choropleth 1: murder rate (sequential scale).}

Create a choropleth mapping \texttt{Murder} (rate per 100,000) to fill color.

You must:
\begin{itemize}
  \item Join \texttt{state\_data} to \texttt{us} on \texttt{region}.
  \item Use a \textbf{sequential} color scale (e.g., \texttt{scale\_fill\_viridis\_c()}).
  \item Use \texttt{coord\_quickmap()} and \texttt{theme\_void()}.
  \item Include a clear legend title.
\end{itemize}

\textbf{Pseudo-code structure:}

\begin{verbatim}
us_merged <- us %>% left_join(________, by = "________")

ggplot(us_merged, aes(x = long, y = lat, group = group,
                       fill = ________)) +
  geom_polygon(color = "white", linewidth = 0.3) +
  coord_quickmap() +
  scale_fill_viridis_c(name = "________") +
  theme_void()
\end{verbatim}

Save as:
\begin{itemize}
  \item \texttt{figures/murder\_rate\_map.png}
\end{itemize}

\item \textbf{Choropleth 2: life expectancy (diverging scale).}

Create a choropleth of \texttt{Life.Exp} using a \textbf{diverging} color scale centered on the national mean.

\textbf{Pseudo-code structure:}

\begin{verbatim}
state_data <- state_data %>%
  mutate(life_dev = Life.Exp - mean(Life.Exp))

us_merged <- us %>% left_join(________, by = "________")

ggplot(us_merged, aes(x = long, y = lat, group = group,
                       fill = ________)) +
  geom_polygon(color = "white", linewidth = 0.3) +
  coord_quickmap() +
  scale_fill_gradient2(low = "________", mid = "________",
                       high = "________", midpoint = 0,
                       name = "________") +
  theme_void()
\end{verbatim}

Save as:
\begin{itemize}
  \item \texttt{figures/life\_exp\_diverging.png}
\end{itemize}

\item \textbf{Choropleth 3: binned variable (categorical scale).}

Create a choropleth of \texttt{HS.Grad} (high-school graduation rate) using a \textbf{binned/categorical} color scale.

You must:
\begin{itemize}
  \item Create a new variable that bins \texttt{HS.Grad} into 4 categories using \texttt{cut()}.
  \item Use \texttt{scale\_fill\_brewer()} with an appropriate palette.
  \item Justify your bin boundaries in the write-up.
\end{itemize}

\begin{verbatim}
state_data <- state_data %>%
  mutate(grad_cat = cut(HS.Grad,
    breaks = c(0, ________, ________, ________, 100),
    labels = c("________", "________", "________", "________")))
\end{verbatim}

Save as:
\begin{itemize}
  \item \texttt{figures/hs\_grad\_binned.png}
\end{itemize}

\item \textbf{Map vs dot plot comparison.}

For the same variable (\texttt{Murder}), create:
\begin{itemize}
  \item The choropleth you already made (reuse from Question 2).
  \item A \textbf{dot plot} showing the top 20 states by murder rate.
\end{itemize}

\begin{verbatim}
top_20 <- state_data %>%
  arrange(desc(________)) %>%
  slice_head(n = 20) %>%
  mutate(region = reorder(region, ________))

ggplot(top_20, aes(x = ________, y = region)) +
  geom_point(size = 2.5) +
  theme_classic() +
  labs(x = "________", y = NULL)
\end{verbatim}

Save as:
\begin{itemize}
  \item \texttt{figures/murder\_dotplot.png}
\end{itemize}

\item \textbf{Interpretation (write-up required).}

Write 12--16 sentences addressing:
\begin{itemize}
  \item What spatial pattern does the murder-rate choropleth suggest? Is there a regional clustering?
  \item For life expectancy, does the diverging scale help you see above/below-average states more clearly than a sequential scale would?
  \item What did your binning choices for HS graduation rate reveal or hide?
  \item Compare the choropleth vs the dot plot for murder rate: when would you prefer each?
  \item Name one concrete plotting choice you made (projection, color scale, binning) and why it helped interpretability.
\end{itemize}

\item \textbf{Submit your work (GitHub or Canvas; proof required).}

\begin{itemize}
  \item \textbf{GitHub option:}
\begin{verbatim}
git status
git add .
git commit -m "Geospatial lab: choropleths + state data"
git push
\end{verbatim}

  \item \textbf{Canvas option:}
  Upload required files to Canvas.
\end{itemize}

Include proof in write-up:
\begin{itemize}
  \item GitHub: output of \texttt{git status} and \texttt{git log -1}
  \item Canvas: confirmation note + list of uploaded files
\end{itemize}

\end{enumerate}

\section*{Optional Extension: USArrests Choropleth}

Using the \texttt{USArrests} dataset, create a choropleth of \texttt{Assault} arrests per 100,000.
\begin{itemize}
  \item Join \texttt{USArrests} to map boundaries (you will need to add a \texttt{region} column from \texttt{tolower(rownames(USArrests))}).
  \item Compare a continuous vs binned version.
\end{itemize}

\begin{verbatim}
arrests <- USArrests %>%
  mutate(region = tolower(rownames(USArrests)))
us_arrests <- us %>% left_join(arrests, by = "region")

ggplot(us_arrests, aes(x = long, y = lat, group = group,
                        fill = Assault)) +
  geom_polygon(color = "white") +
  coord_quickmap() +
  scale_fill_viridis_c() +
  theme_void()
\end{verbatim}

\section*{Checklist}
\begin{itemize}
  \item Script runs top-to-bottom
  \item Required figures saved
  \item Write-up complete
  \item Submitted via GitHub or Canvas
\end{itemize}

\end{document}
